\documentclass[11pt]{article}
\usepackage[a4paper,margin=25mm]{geometry}
\usepackage[T1]{fontenc}
\usepackage{lmodern,microtype}
\usepackage{amsmath,amssymb,amsthm,mathtools}
\usepackage{booktabs,array}
\usepackage[hidelinks]{hyperref}
\usepackage{fancyhdr}
\hypersetup{pdftitle={Two-colour rigidity and boundary obstructions for entropy extremizers},pdfauthor={B2 research continuation},pdfsubject={Further results for Q4296 and Q4660},pdfkeywords={midpoint colouring, affine cap, entropy, support perturbation, exact certificate}}
\setlength{\parskip}{3pt}
\setlength{\emergencystretch}{2em}
\setlength{\parindent}{0pt}
\setlength{\headheight}{14pt}
\pagestyle{fancy}\fancyhf{}
\fancyhead[L]{\small B2 research continuation II}
\fancyhead[R]{\small 13 October 2026}
\fancyfoot[C]{\thepage}
\newtheorem{theorem}{Theorem}[section]
\newtheorem{lemma}[theorem]{Lemma}
\newtheorem{proposition}[theorem]{Proposition}
\newtheorem{corollary}[theorem]{Corollary}
\theoremstyle{remark}
\newtheorem{remark}[theorem]{Remark}
\newcommand{\F}{\mathbb F_3}
\newcommand{\Fp}{\mathbb F_p}
\newcommand{\Prob}{\mathbb P}
\newcommand{\supp}{\operatorname{supp}}
\newcommand{\aff}{\operatorname{aff}}
\newcommand{\codim}{\operatorname{codim}}
\newcommand{\Dplus}{\Delta_+}
\newcommand{\Dminus}{\Delta_-}
\title{\vspace{-12mm}\textbf{Two-colour rigidity and boundary obstructions\\for entropy extremizers}\\[4pt]\large Further results for the B2 midpoint and entropy programs}
\author{Research continuation from the supplied handoff}
\date{13 October 2026}
\begin{document}
\maketitle\vspace{-6mm}
\begin{abstract}
An admissible two-colour partition of $\F^n$ has a line-free remainder only
when that remainder is empty, a singleton, two points, or a four-point affine
parallelogram. An elementary plane-incidence proof establishes this
classification in every dimension. A full-fibre inheritance lemma then shows
that, in any fixed direction, the number of parallel maximal affine subspaces
in a nonempty two-colour remainder is $1$, $2$, or $4$, with parallelogram
structure in the last case. Colour-reusing ambient extension gives density
$1/10$ examples whose largest affine subspace has dimension exactly $n-k-2$.
For the sum--difference entropy ratio, a support-boundary expansion gives a
finite test for leading-order support enlargement. In particular, no finite
law of ratio greater than one that fully occupies an arithmetic progression
can be globally extremal. Exact rational arithmetic also certifies the
improved bound $C_*>1.4051810329$. These results do not remove the general
$O(k\log k)$ codimension loss or settle whether $C_*=2$.
\end{abstract}

\section{Scope and relation to the previous results}

The two targets and their stable statement identifiers are
\begin{center}
\begin{tabular}{ll}
Q4296 & \texttt{67d9a47fe0b199024fdb}\\
Q4660 & \texttt{4e69214fdfaf0bc4e852}.
\end{tabular}
\end{center}
The previous continuation \cite{previous} proved an all-dimension midpoint
bound of codimension
\[
 O_p\!\left(k[\log(2k)+\log(1/\alpha)]\right)
\]
using the established bounded-torsion PFR theorem \cite{pfr}, and certified an entropy ratio of
$1.4051809660756335725\ldots$. The present note adds structural theorems,
exact small-parameter values, a stronger explicit obstruction, and a support
perturbation method. The primary midpoint paper gives the quadratic bound
and asks for the exact extremal function \cite{midpoint}.

\textbf{Source-location correction.} The catalog attaches ``Conjecture 1.10''
to Q4660. In the primary source \cite{li}, Conjecture 1.10 is the separate
sum--product conjecture. The sum--difference sharpness question stated in
Q4660 occurs in Section 8. This note addresses Q4660's mathematical statement,
not Conjecture 1.10. The original catalog record is preserved unmodified;
the correction is recorded separately in the accompanying provenance.

\textbf{Status.} The general statements below have complete written proofs.
The numerical lower bound has an exact rational certificate. The programs
provide independent finite controls, not a formal verification of the
all-dimension arguments. No external peer review, proof-assistant
verification, or worldwide-priority claim is made. Neither complete catalog
target is marked solved. The other targets and the earlier attempt counters
are unchanged.

\section{Exact classification of two-colour cap remainders}

A partition
\[
 \F^n=A\sqcup B\sqcup R
\]
is \emph{admissible} when two points in either one of the colour classes
$A,B$ never have their midpoint in $R$. Empty colour classes are permitted.
No density assumption is imposed in this section. A \emph{cap} is a set
containing no affine line; over $\F$, a line has three points. An affine
parallelogram means a set $a+\{0,u,v,u+v\}$ with $u,v$ linearly independent.

\begin{theorem}[Two-colour cap classification]\label{thm:cap}
For every $n\geq0$, if an admissible two-colour remainder $R\subseteq\F^n$
is a cap, then $R$ is empty, a singleton, two distinct points, or an affine
parallelogram. Conversely, each of these configurations is the remainder
of an admissible partition in every ambient space containing it.
In particular, $|R|\leq4$ and a nonempty $R$ has affine dimension at most two.
\end{theorem}

\begin{lemma}[No three-point plane section]\label{lem:three}
Under the hypotheses of Theorem~\ref{thm:cap}, no affine plane meets $R$
in exactly three points.
\end{lemma}
\begin{proof}
Such three points would be noncollinear. Affine coordinates on their plane
therefore put them at $0,e_1,e_2$. The three points
\[
 e_1+e_2,\qquad e_1-e_2,\qquad -e_1+e_2
\]
are outside $R$, and their pairwise midpoints are $e_1,e_2,0$.
They must be pairwise differently coloured, which is impossible with two
colours. The argument only uses points in the plane, so it applies to every
plane section of the full ambient partition.
\end{proof}

\begin{lemma}[Caps in an affine plane]\label{lem:planecap}
A cap in $\F^2$ has at most four points. Every four-point cap is an affine
parallelogram.
\end{lemma}
\begin{proof}
Fixing a point of a nonempty cap, the four lines through it contain at most
one other cap point each. The cap therefore has size at most five.
If it had size five, this argument at each cap point would show that no line
meets the cap in exactly one point. But any parallel class of three lines
partitions the five points into parts of size at most two, necessarily of
sizes $2,2,1$. This is a contradiction.

For a four-point cap, put three of its points at $0,e_1,e_2$. The fourth
point can only be $e_1+e_2$, $e_1-e_2$, or $-e_1+e_2$: the other unused
points complete a line through two of the first three. Each of the three
possibilities is a parallelogram, possibly with a different choice of
opposite vertices.
\end{proof}

\begin{proof}[Proof of Theorem~\ref{thm:cap}: the dimension bound]
Let $W$ be any affine three-space and set $S=R\cap W$, $s=|S|$.
By Lemmas~\ref{lem:three} and \ref{lem:planecap}, every affine plane of $W$
meets $S$ in $0$, $1$, $2$, or $4$ points. Write $n_j$ for the number of
planes with intersection size $j$.

There are $39$ affine planes in $\F^3$: the nonzero linear functionals,
up to their two nonzero scalar multiples, give $13$ plane directions,
each with three translates. A point lies on $13$ planes. A pair lies on
four planes, because the linear functionals vanishing on its one-dimensional
difference space form a two-dimensional space, with four projective classes.
Every triple in the cap is noncollinear and determines one plane. Counting
planes, point incidences, pair incidences, and triple incidences gives
\begin{align}
 n_0+n_1+n_2+n_4&=39,\label{eq:mom0}\\
 n_1+2n_2+4n_4&=13s,\label{eq:mom1}\\
 n_2+6n_4&=4\binom{s}{2},\label{eq:mom2}\\
 4n_4&=\binom{s}{3}.\label{eq:mom3}
\end{align}
Consequently,
\begin{align*}
 n_4&=\frac14\binom{s}{3},&
 n_2&=4\binom{s}{2}-\frac32\binom{s}{3},\\
 n_1&=13s-8\binom{s}{2}+2\binom{s}{3},&
 n_0&=39-13s+4\binom{s}{2}-\frac34\binom{s}{3}.
\end{align*}
Three parallel planes cover $W$, and each contains at most four cap points,
so $s\leq12$. For $s=3$, equation~\eqref{eq:mom3} is already impossible.
For $5\leq s\leq12$, the following table rules out nonnegative integral
values of all the counts.
\begin{center}
\begin{tabular}{rrrrr}
\toprule
$s$ & $n_0$ & $n_1$ & $n_2$ & $n_4$\\
\midrule
5 & $13/2$ & 5 & 25 & $5/2$\\
6 & 6 & $-2$ & 30 & 5\\
7 & $23/4$ & $-7$ & $63/2$ & $35/4$\\
8 & 5 & $-8$ & 28 & 14\\
9 & 3 & $-3$ & 18 & 21\\
10 & $-1$ & 10 & 0 & 30\\
11 & $-31/4$ & 33 & $-55/2$ & $165/4$\\
12 & $-18$ & 68 & $-66$ & 55\\
\bottomrule
\end{tabular}
\end{center}
The only possible sizes are therefore $0,1,2,4$. If $s=4$, then $n_4=1$,
so the four points lie in one affine plane.

If $R$ contained four affinely independent points, their affine span would
be a three-space $W$. Its section $S$ would have at least four points and
would not be contained in a plane, contradicting the preceding conclusion.
Thus $\dim\aff(R)\leq2$ when $R$ is nonempty. The plane lemmas now give
the claimed list, including the exclusion of three points. For ambient
dimension less than three the same plane lemmas apply directly.
\end{proof}

\subsection{Extension without spending another colour}

The converse and the later constructions use an elementary extension
principle that is useful independently of the classification.

\begin{lemma}[Colour-reusing ambient extension]\label{lem:extension}
Let $p$ be an odd prime and $k\geq2$. An admissible $k$-colour partition
of an affine subspace $U\subseteq\Fp^n$ extends to an admissible partition
of $\Fp^n$ with the same remainder and no additional colour labels.
\end{lemma}
\begin{proof}
Translate $U$ to a linear subspace. In $\Fp^n/U$, pair each nonzero quotient
point $q$ with $-q$ and arbitrarily orient each pair. Colour every point of
the positively oriented fibre with colour $1$, and every point of the
opposite fibre with colour $2$. Retain the original partition on $U$.

Two outside points can have midpoint in $U$ only if their nonzero quotient
points are opposites, in which case the colours differ. A pair with exactly
one endpoint in $U$ has midpoint outside $U$. Pairs wholly in $U$ obey the
original condition. Since the remainder is unchanged and contained in $U$,
these observations prove admissibility.
\end{proof}

\begin{proof}[Completion of Theorem~\ref{thm:cap}: realizability]
The empty remainder is immediate. A singleton is the remainder of the
zero-dimensional partition and extends by Lemma~\ref{lem:extension}.
Two points are realized on their three-point affine line by colouring the
third point, and then extending. For the square
$Q=\{1,2\}\times\{1,2\}\subseteq\F^2$, colour points with first coordinate
zero by colour $1$, and points with second coordinate zero and first
coordinate nonzero by colour $2$. A monochromatic midpoint has a zero
coordinate and cannot belong to $Q$. Affine invariance and extension give
every required parallelogram in every containing ambient space.
\end{proof}

\section{Parallel maximal subspaces and extremal consequences}

\begin{lemma}[Inheritance by full fibres]\label{lem:fibre}
Let $G=\Fp^n=A_1\sqcup\cdots\sqcup A_k\sqcup R$ be admissible, with $p$
odd, and let $\pi:G\to G/V$ be a linear quotient. Set
\[
 R_V=\{q\in G/V:\pi^{-1}(q)\subseteq R\}.
\]
Then $R_V$ is an admissible $k$-colour remainder in the quotient. For each
$q\notin R_V$, any chosen point of $\pi^{-1}(q)\setminus R$ can supply
$q$'s colour.
\end{lemma}
\begin{proof}
Suppose two chosen representatives $x,y$ have the same colour. If the
midpoint of their quotient points belonged to $R_V$, its entire fibre
would lie in $R$, including $(x+y)/2$. This contradicts admissibility in
$G$. The proof does not require a partially occupied fibre to be
monochromatic and does not assert preservation of remainder density.
\end{proof}

\begin{corollary}[Rigidity of parallel maximal flats]\label{cor:parallel}
Suppose an admissible two-colour remainder $R\subseteq\F^n$ contains no
affine $(d+1)$-subspace. For every $d$-dimensional linear subspace $V$, the
number of complete cosets of $V$ contained in $R$ is $0,1,2$, or $4$.
In the four-coset case, their quotient points form an affine parallelogram.
In particular, if $V$ is the direction of a largest affine subspace of a
nonempty $R$, the number is $1,2$, or $4$, and never three.
\end{corollary}
\begin{proof}
An affine line in $R_V$ would lift to an affine $(d+1)$-subspace wholly
inside $R$. Thus $R_V$ is a cap. Apply Lemma~\ref{lem:fibre} and
Theorem~\ref{thm:cap}.
\end{proof}
The four-coset case is sharp in every $d$: take
$R=Q\times\F^d$, with $Q$ the four-point cap, and lift its two-colour
partition by Cartesian product. Projection of any affine subspace of $R$
into $Q$ must be a point, since a positive-dimensional affine image would
contain a line. The largest affine dimension is therefore exactly $d$.
Ambient extension is again allowed.

\begin{corollary}[Exact first two-colour values]\label{cor:small}
Let $f(n,2)$ be the minimum, over admissible two-colour partitions with
$|R|\geq3^n/10$, of the largest affine dimension in $R$. Then
\[
 f(1,2)=f(2,2)=f(3,2)=0,\qquad f(4,2)=1.
\]
For every $n\geq4$, a remainder of density at least $1/10$ contains a line.
In dimension three, every line-free remainder meeting the density threshold
is a four-point parallelogram.
\end{corollary}
\begin{proof}
Singletons give the first two examples: their densities are $1/3$ and $1/9$.
A four-point cap embedded in $\F^3$ has density $4/27$ and largest affine
dimension zero. In dimension four the density condition requires at least
nine points, exceeding the cap bound of four, and hence forces a line.
The example $Q\times\F\times\{0\}\subseteq\F^4$, extended from its
three-dimensional affine span, has density $12/81$ and largest affine
dimension one. The same cap bound forces a line in all higher dimensions.
\end{proof}

\begin{proposition}[A $k+2$ codimension obstruction]\label{prop:block}
Let $k\geq2$, $b=\lceil\log_3(10k)\rceil$, and $n\geq kb+2$.
There is an admissible $k$-colour partition of $\F^n$ with density
$|R|/3^n\geq1/10$ for which the largest affine subspace in $R$ has dimension
exactly $n-k-2$. In particular,
\[
 f(n,k)\leq n-k-2\qquad
 \text{for }n\geq k\lceil\log_3(10k)\rceil+2.
\]
\end{proposition}
\begin{proof}
Write coordinates as $(z_1,\ldots,z_k,t,y)$, with
$z_i\in\F^b$, $t\in\F^2$, and $y\in\F^{n-kb-2}$. Define
\[
 R=\{(z_1,\ldots,z_k,t,y):z_i\neq0\ \text{for all }i,\ t=0\}.
\]
Inside $t=0$, colour every point outside $R$ by its first zero block.
A pair of colour $i$ has a zero $i$th block at its midpoint, so cannot have
midpoint in $R$. Extend to the whole ambient space by
Lemma~\ref{lem:extension}, reusing colours $1,2$. Explicitly, for $t\neq0$
use its first nonzero coordinate as colour $1$ or $2$.

The density is
\[
 \frac{|R|}{3^n}=\frac19(1-3^{-b})^k
 \geq\frac19(1-k3^{-b})\geq\frac1{10}.
\]
Every affine subspace $L\subseteq R$ satisfies the two independent
constraints $t=0$. Its image in block $i$ is an affine space $a_i+W_i$
avoiding zero, so $a_i\notin W_i$. Some linear functional $\ell_i$ has
$\ell_i(W_i)=0$ and $\ell_i(a_i)=1$. The $k$ resulting block constraints
are independent of one another and of the two $t$-constraints, because
their coordinate supports are disjoint. Hence $\dim L\leq n-k-2$.
Equality is attained by setting the first coordinate of each $z_i$ to one
and $t=0$, while leaving all other coordinates free.
\end{proof}

\begin{remark}[What these conclusions do not imply]
The proposition strengthens the previous $k$-codimension construction by
two and rules out a universal exact guarantee of dimension $n-k$. It does
not change the order of the lower obstruction or improve the general
$O(k\log k)$ upper codimension bound from the previous continuation.
Full-fibre inheritance does not preserve density and cannot by itself be
iterated into a linear-codimension theorem.
\end{remark}

\section{A support-boundary calculus for the entropy ratio}

Let $P$ be a nondegenerate probability law on a finite set $A\subseteq\mathbb R$,
and let $X,X'$ be independent with law $P$. Use natural logarithms and write
\[
 \Dplus(P)=H(X+X')-H(X),\qquad
 \Dminus(P)=H(X-X')-H(X),\qquad
 \rho(P)=\frac{\Dminus(P)}{\Dplus(P)}.
\]
The denominator is positive. Indeed, the law of $X+X'$ is a nontrivial
mixture of distinct translates of $P$, all with entropy $H(P)$, and entropy
is strictly concave. The translates are distinct because a finite
nonempty support cannot be invariant under a nonzero real translation.
The ratio does not depend on the logarithm base.

Write $C_*$ for the least universal sum--difference constant in Q4660.
The bound $C_*\leq2$ is established in the literature; sharpness is the
question recorded in Section 8 of \cite{li}. The dimension equivalence in
\cite{makkuva} is also relevant: independent tensor products preserve a
ratio rather than increasing it. No new tensorization theorem is claimed
here.

\begin{theorem}[Support-boundary expansion]\label{thm:boundary}
For $z\notin A$, put $P_\varepsilon=(1-\varepsilon)P+\varepsilon\delta_z$,
and define the two escape masses
\[
 u_+(z)=\Prob(z+X\notin A+A),\qquad
 u_-(z)=\Prob(z-X\notin A-A).
\]
As $\varepsilon\downarrow0$,
\begin{equation}\label{eq:boundary}
 \rho(P_\varepsilon)=\rho(P)+
 \frac{L_P(z)}{\Dplus(P)}\varepsilon\log(1/\varepsilon)
 +O_{P,z}(\varepsilon),
\end{equation}
where
\begin{equation}\label{eq:L}
 L_P(z)=2u_-(z)-1-\rho(P)\bigl(2u_+(z)-1\bigr).
\end{equation}
In particular, $L_P(z)>0$ implies strict ascent for every sufficiently
small positive $\varepsilon$. A finitely supported global maximizer must
have $L_P(z)\leq0$ for every $z\notin A$.
\end{theorem}
\begin{proof}
Adding a new input atom gives
\[
 H(P_\varepsilon)=H(P)+\varepsilon\log(1/\varepsilon)+O_P(\varepsilon).
\]
For a finite output law depending polynomially on $\varepsilon$, an atom
whose mass is positive at $\varepsilon=0$ contributes a change of order
$O(\varepsilon)$ to entropy. A previously zero atom whose mass is
$c\varepsilon+O(\varepsilon^2)$, $c>0$, contributes
$c\varepsilon\log(1/\varepsilon)+O(\varepsilon)$. An atom first appearing
at order $\varepsilon^2$ contributes only
$O(\varepsilon^2\log(1/\varepsilon))=o(\varepsilon)$.

In the sum law, the two input cross terms have total first-order mass
$2\varepsilon$, and the part landing outside $A+A$ is
$2u_+(z)\varepsilon$. In the difference law the cross terms are supported
on $z-A$ and $A-z$. Symmetry of $A-A$ shows that each has escape mass
$u_-(z)\varepsilon$. Coincident output atoms are aggregated; no
disjointness or independence between these escape events is assumed.
Thus
\begin{align*}
 \Dplus(P_\varepsilon)&=\Dplus(P)
 +(2u_+(z)-1)\varepsilon\log(1/\varepsilon)+O_{P,z}(\varepsilon),\\
 \Dminus(P_\varepsilon)&=\Dminus(P)
 +(2u_-(z)-1)\varepsilon\log(1/\varepsilon)+O_{P,z}(\varepsilon).
\end{align*}
Divide the two expansions. Since $\Dplus(P)>0$ and
$(\varepsilon\log(1/\varepsilon))^2=o(\varepsilon)$, the quotient error
remains $O_{P,z}(\varepsilon)$, proving~\eqref{eq:boundary}. Its positive
leading term dominates that error when $L_P(z)>0$.
\end{proof}

\begin{corollary}[A finite boundary test]\label{cor:finite}
Suppose $\rho(P)>1$. Any atom $z$ with $L_P(z)>0$ belongs to
\[
 (A+A-A)\setminus A.
\]
Consequently, testing all leading-order support-enlargement directions
requires at most $|A|^3$ candidate locations. Passing the test is a
necessary condition for global optimality, not a sufficient condition.
\end{corollary}
\begin{proof}
If $z\notin A+A-A$, then $z+x\notin A+A$ and $z-x\notin A-A$ for
every $x\in A$. Thus $u_+(z)=u_-(z)=1$ and
$L_P(z)=1-\rho(P)<0$. The displayed finite set includes every other
possible new atom with nonnegative leading coefficient.
\end{proof}

\begin{corollary}[No finite full-AP global maximizer]\label{cor:AP}
Suppose the support of $P$ is a nontrivial finite arithmetic progression,
with positive probability at every one of its points, and $\rho(P)>1$.
Then $P$ is not a global maximizer of the sum--difference ratio. One of
the two adjacent endpoint atoms gives strict ascent at all sufficiently
small positive masses.
\end{corollary}
\begin{proof}
Affine changes of the random variable preserve the ratio. Normalize the
support to $A=\{0,\ldots,N\}$, with probabilities $p_j>0$. For a new
right-hand atom $z=N+1$, the only input producing a sum outside $A+A$ is
$X=N$, and the only input producing a difference outside $A-A$ is $X=0$.
Therefore
\[
 L_R=2p_0-1+\rho(1-2p_N).
\]
At the left-hand atom $z=-1$ the roles reverse:
\[
 L_L=2p_N-1+\rho(1-2p_0).
\]
Their sum is
\begin{equation}\label{eq:endpointsum}
 L_R+L_L=2(\rho-1)(1-p_0-p_N).
\end{equation}
For $N\geq2$, the right side is strictly positive. If $N=1$, then
$L_R=(1+\rho)(p_0-p_1)$ and $L_L=-L_R$. Equality of the endpoint masses
would give the symmetric two-point law with $\rho=1$, contrary to the
hypothesis. Hence at least one endpoint has positive coefficient in every
case. Apply Theorem~\ref{thm:boundary}.
\end{proof}
The same argument can be repeated after each addition: every finite
full-AP law of ratio above one admits an infinite chain of strictly
improving finite full-AP extensions. If the initial probabilities are
rational, the added masses can all be chosen rational. This does not
identify the limiting ratio or assert that it approaches two.

\begin{corollary}[Heavy endpoint obstruction]\label{cor:heavy}
A finitely supported real law with ratio greater than one and probability
at least $1/2$ at its minimum or its maximum cannot be globally extremal.
\end{corollary}
\begin{proof}
By reflection consider $a=\min A$, $b=\max A$, and $p_a\geq1/2$.
At $z=2b-a>b$, the only old sum reached by $z+X$ occurs at $X=a$, and
the only old difference reached by $z-X$ occurs at $X=b$. Hence
\[
 u_+(z)=1-p_a,\qquad u_-(z)=1-p_b,
 \qquad L_P(z)=1-2p_b-\rho(1-2p_a).
\]
This is positive: $p_b\leq1-p_a\leq1/2$, and equality in both terms
could only be the fair two-point law, whose ratio is one. Apply the
support-boundary theorem.
\end{proof}

\begin{remark}[The boundary test is not an optimality certificate]
For the uniform law on $A=\{0,1,4\}$,
\[
 \rho=\frac{(2/3)\log3}{\log3-(2/3)\log2}>1.
\]
The candidate set $(A+A-A)\setminus A$ is
$\{-4,-3,-2,-1,2,3,5,7,8\}$. At each candidate, the old-support capture
mass $1-u_+$ is $1/3$, while $1-u_-$ is either $1/3$ or $2/3$.
The corresponding coefficients are $(1-\rho)/3$ or $-(1+\rho)/3$,
both strictly negative; all other atoms also have negative coefficient.
This law nevertheless has ratio less than $6/5$, well below the witness
in the next section. Thus failure to find a positive leading coefficient
does not prove global optimality, and the preceding corollaries do not
exclude every finite support. The bound $\rho<6/5$ follows from
$8\log3>12\log2$, equivalently $3^2>2^3$.
\end{remark}

\section{A stronger exact entropy certificate}

\begin{proposition}[Twelve-point witness]\label{prop:witness}
Let $X$ have support $\{0,1,\ldots,11\}$ and probability numerators
\[
\begin{array}{rrrr}
88162313713840,&10154514342618,&1485728499913,&177389643984,\\
18240191721,&1665031585,&137417507,&10383619,\\
725058,&47123,&2867,&165
\end{array}
\]
with common denominator $10^{14}$. Then
\begin{gather*}
1.4051810329220149462009831947934829507300<\rho(X),\\
\rho(X)<1.4051810329220149462009831947934829507301.
\end{gather*}
In particular,
\[
 C_*>\frac{14051810329}{10^{10}}=1.4051810329.
\]
\end{proposition}

\begin{proof}[Exact certificate]
For integer input weights $w_i$ of total $D$, compute output weights by
integer convolution,
\[
 s_t=\sum_{i+j=t}w_iw_j,\qquad d_t=\sum_{i-j=t}w_iw_j.
\]
Their common total is $D^2$. For a nonnegative count vector $c$ of positive
total $N$, entropy is
\[
 H(c)=\log N-\frac1N\sum_{c_i>0}c_i\log c_i.
\]
The included program \texttt{certify\_entropy.py} uses only exact integers
and fractions. It encloses each logarithm by range reduction to $r\in[1,2]$,
followed by $t=(r-1)/(r+1)\in[0,1/3]$ and
\[
 L_J(r)=2\sum_{j=0}^{J-1}\frac{t^{2j+1}}{2j+1},\qquad
 0\leq\log r-L_J(r)\leq
 \frac{2t^{2J+1}}{(2J+1)(1-t^2)}.
\]
It takes $J=90$ and rounds the resulting rational bounds outward to the
$10^{-65}$ grid before subsequent entropy arithmetic. It proves positivity
of both increments, bounds their quotient with the correct monotone
endpoints, and directly proves
\[
 \underline{\Dminus}
 -\frac{14051810329}{10^{10}}\,\overline{\Dplus}>0.
\]
The file \texttt{ENTROPY-CERTIFICATE.json} records all convolution weights,
interval endpoints, and the positive-gap certificate. The unrounded
rational bounds lie strictly inside the displayed decimal endpoints,
which are outward-rounded to forty decimal places. This finite rational
calculation proves the stated inequalities.
\end{proof}

The improvement over the preceding eight-point certificate is approximately
$6.68464\cdot10^{-8}$ in the ratio. It is a numerical increment, not a
large change in the known universal range. The more substantial entropy
conclusion is Corollary~\ref{cor:AP}: the displayed law itself, and every
other finite full-AP candidate with ratio above one, still admits a
strictly improving extension.

As a concrete check on the perturbation mechanism, start with weights
$(8830,1009,143,16,2)$ of total $10^4$, multiply each weight by $999999$,
and append a sixth weight $10000$. This realizes
$\varepsilon=10^{-6}$ at the next endpoint. Exact certification gives
\[
 1.405025826013192214095918\ldots
 \quad\longrightarrow\quad
 1.405056049166641606653721\ldots.
\]
The complete before-and-after certificates appear in
\texttt{BOUNDARY-CHECKS.json}; this is a demonstrative ascent, not the
strongest witness in the package.

\section{Verification, negative controls, and remaining gaps}

The three proof-control programs require only Python's standard library:
\begin{verbatim}
python certify_entropy.py
python verify_midpoint.py
python verify_entropy_boundary.py
\end{verbatim}
The entropy verifier recomputes the exact twelve-point certificate and
checks the previous eight-point law as a baseline control. The midpoint
verifier enumerates all $512$ plane subsets, confirming cap counts
$1,9,36,72,54$ at sizes $0,1,2,3,4$ and admissible cap counts
$1,9,36,54$ at sizes $0,1,2,4$. It independently constructs the $39$
planes of $\F^3$, checks point, pair, and noncollinear-triple incidences,
and evaluates every possible moment row for $0\leq s\leq12$.
It also examines all $19683$ two-colour labelled plane partitions,
finds $1371$ admissible ones, and checks all $23196$ inherited quotient
colour assignments generated by arbitrary allowed representative colours.
The entropy-boundary controls check $2238$ new-output mass identities,
$16$ AP endpoint identities, $120$ heavy-endpoint instances, and the exact
finite ascent above. Complete boundary reports find eleven positive leading
directions among the twenty-two candidates for the twelve-point law, and
strictly negative coefficients at every candidate for the uniform
$\{0,1,4\}$ negative control.

The package preserves a negative control against an overstrong midpoint
strategy. In $\F^3$, let
\begin{align*}
 A=\{&(0,1,2),(0,2,2),(1,1,0),\\
     &(1,2,1),(2,0,0),(2,0,1)\},\qquad
 B=\{(x,y,z):y+z=2\},
\end{align*}
and let $R$ be the complement. This is an admissible partition; both $A$
and $R$ affinely span the ambient space, while $R$ contains no affine plane.
The verifier checks these statements exactly. In particular, a removal
argument that assumes every colour has proper affine span whenever the
remainder spans the ambient space is invalid.

Floating-point stationary searches on longer consecutive supports,
correlated grids, and other support patterns are retained separately in
\texttt{exploratory/}. They suggested the rational witness but supply no
global upper bound. MILP outputs in that directory are exploratory solver
records, not certificates of dimension-free statements. The elementary
classification proof does not depend on the search that first suggested it.

\textbf{Remaining gaps.} The exact general function $f(n,k)$ is not
determined. In particular, the prior $O(k\log(2k))$ upper codimension bound
has not been improved to $O(k)$ here. The full-fibre result counts parallel
complete subspaces but does not control partially occupied fibres.
For entropy, the universal upper bound remains two. Finite-support
boundary tests impose necessary conditions and eliminate full-AP and
heavy-endpoint candidates; they neither certify a global optimum nor
exclude all finite supports. There is no claim of a solution to the
separate entropic sum--product conjecture.

\begingroup\small
\begin{thebibliography}{9}\setlength{\itemsep}{3pt}\setlength{\parskip}{0pt}
\bibitem{midpoint}
Tomasz Ko\'sciuszko, \emph{Arithmetic Progressions in Midpoint Colourings},
arXiv:2610.11847v1 (8 October 2026), Problem 4 and Theorem 3.
\url{https://arxiv.org/html/2610.11847v1}.
\bibitem{li}
Rupert Li, \emph{The Entropic Sum-Product Phenomenon},
arXiv:2607.29042v1 (2026), Conjecture 1.10 and Section 8.
\url{https://arxiv.org/html/2607.29042v1}.
\bibitem{makkuva}
Ashok Vardhan Makkuva and Yihong Wu,
\emph{Equivalence of additive-combinatorial linear inequalities for Shannon
entropy and differential entropy}, arXiv:1601.07498v2, Proposition 1.
\url{https://arxiv.org/html/1601.07498v2}.
\bibitem{pfr}
W. T. Gowers, Ben Green, Freddie Manners, and Terence Tao,
\emph{Marton's conjecture in abelian groups with bounded torsion},
arXiv:2404.02244v2, Theorem 1.1.
\url{https://arxiv.org/html/2404.02244v2}.
\bibitem{previous}
\emph{Logarithmic codimension bounds for midpoint colourings and certified
improvements for the entropic sum--difference problem}, previous B2
continuation (13 October 2026), supplied as
\texttt{B2-new-results-2026-10-13.zip}. This is a local research manuscript,
not a claim of independent publication.
\end{thebibliography}
\endgroup
\end{document}
